% 原书第六篇6.5节227--237；印刷258--259，PDF267--268；
% 书末答案348--349，PDF357--358。
\radioactivityproblem{227}{258}
一制剂由两种半衰期不同的 $\beta$ 放射性核素组成。
按时间 $t$（小时）测得
\[
\begin{array}{c|rrrrrrrrr}
t&0&1&2&3&5&7&10&14&20\\\hline
\ln A&4.10&3.60&3.10&2.60&2.06&1.82&1.60&1.32&0.90
\end{array}
\]
估算两种核素的半衰期及 $t=0$ 时核数比。
\exsolution
总活度满足双指数模型
\[
 A(t)=A_{f0}e^{-\lambda_f t}
 +A_{s0}e^{-\lambda_s t},\quad\lambda_f>\lambda_s.
\]
末期 $t=14,20$ 两个点给慢分量的近似斜率
$\lambda_s\simeq(1.32-0.90)/6=0.070\,\mathrm{h^{-1}}$，
于是 $A_{s0}\simeq e^{0.90+20\lambda_s}\simeq10$。
初始总活度为 $e^{4.10}\simeq60.34$，
所以 $A_{f0}\simeq50.34$。
从早期点减去慢分量再拟合快指数，得
$\lambda_f\simeq0.6$--$0.7\,\mathrm{h^{-1}}$。
若对全部九个 $\ln A$ 数据作双指数最小二乘拟合，
一组近似参数为
$\lambda_f\simeq0.701\,\mathrm{h^{-1}}$、
$\lambda_s\simeq0.0640\,\mathrm{h^{-1}}$，
$A_{f0}\simeq53.8$、$A_{s0}\simeq9.01$。
于是
\[
 \boxed{T_f\simeq0.99\,\mathrm h,\qquad
 T_s\simeq10.8\,\mathrm h}.
\]
由于 $N=A/\lambda$，慢/快初始核数之比
\[
 \boxed{\frac{N_{s0}}{N_{f0}}=
 \frac{A_{s0}\lambda_f}{A_{f0}\lambda_s}
 \simeq1.83}.
\]
\excheck 书末印 $T_f=1.6\,\mathrm h$、
$T_s=9.8\,\mathrm h$、核数比10。
将该组合代回\emph{原扫描表中全部九个}\,$\ln A$
不能同时拟合；本稿给出可重复双指数拟合及参数，
并把原书答案与数据矛盾列为勘误，待独立复核。

\radioactivityproblem{228}{258}
反应堆以恒定速率
$q=2.7\times10^9\,\mathrm{s^{-1}}$
生产 $\mathrm{^{32}P}$，
半衰期 $T_{1/2}=14.3\,\mathrm d$。
开始生产后多久该核素的活度达到
$A=1.0\times10^9\,\mathrm{Bq}$？
\exsolution
由生产与衰变的核数方程
$dN/dt=q-\lambda N$，$N(0)=0$，
解得
$N(t)=(q/\lambda)(1-e^{-\lambda t})$，
活度
$A(t)=\lambda N(t)=q(1-e^{-\lambda t})$。
故
\[
 \boxed{t=-\frac1\lambda\ln\left(1-\frac Aq\right)
 =-\frac{T_{1/2}}{\ln2}\ln\left(1-\frac Aq\right)
 \simeq9.5\,\mathrm d}.
\]
\excheck 随时间 $A(t)\to q$，
目标活度必须满足 $A<q$，本题成立。

\radioactivityproblem{229}{258}
一条衰变链 $A_1\xrightarrow{\lambda_1}A_2
\xrightarrow{\lambda_2}\cdots$，
初始只有 $N_{10}$ 个母核 $A_1$，
求 (a) 子核 $A_2$ 随时间的数量；
(b) 子核活度达到最大时刻。
\exsolution
母核
$N_1=N_{10}e^{-\lambda_1t}$，
子核满足
$dN_2/dt=\lambda_1N_1-\lambda_2N_2$，
$N_2(0)=0$。
乘积分因子 $e^{\lambda_2t}$ 积分，得到
\[
 \boxed{N_2(t)=\frac{\lambda_1N_{10}}
 {\lambda_2-\lambda_1}
 (e^{-\lambda_1t}-e^{-\lambda_2t})}.
\]
子核活度 $A_2=\lambda_2N_2$，在
$dN_2/dt=0$ 时极大；
$\lambda_1e^{-\lambda_1t_m}
 =\lambda_2e^{-\lambda_2t_m}$，
所以
\[
 \boxed{t_m=\frac{\ln(\lambda_1/\lambda_2)}
 {\lambda_1-\lambda_2}}.
\]
\excheck 对 $\lambda_1\ne\lambda_2$ 公式成立，
并且 $t_m>0$；两常数相等时取极限见下题。

\radioactivityproblem{230}{258}
若前题两衰变常数相等
$\lambda_1=\lambda_2=\lambda$，
求子核 $A_2$ 数量随时间的函数及达到最大值的时刻。
\exsolution
方程
$dN_2/dt+\lambda N_2=
\lambda N_{10}e^{-\lambda t}$，
乘 $e^{\lambda t}$ 得
$d(N_2e^{\lambda t})/dt=\lambda N_{10}$。
因此
\[
 \boxed{N_2(t)=\lambda N_{10}t e^{-\lambda t}}.
\]
求导
$dN_2/dt=\lambda N_{10}e^{-\lambda t}(1-\lambda t)$，
所以
$\boxed{t_m=1/\lambda=\tau}$。
\excheck 两指数相减型
Bateman公式在 $\lambda_2\to\lambda_1$
时也趋于此式。

\radioactivityproblem{231}{258}
衰变链
$A_1\xrightarrow{\lambda_1}A_2
\xrightarrow{\lambda_2}A_3$
中 $A_3$ 稳定，初始只有 $N_{10}$ 个 $A_1$。
求稳定核 $A_3$ 的累积数目。
\exsolution
根据上题组解，
$N_1=N_{10}e^{-\lambda_1 t}$、
\[
 N_2=\frac{\lambda_1N_{10}}
 {\lambda_2-\lambda_1}
 (e^{-\lambda_1t}-e^{-\lambda_2t}).
\]
整个核链核数守恒，
$N_1+N_2+N_3=N_{10}$，
所以
\[
 \boxed{N_3(t)=N_{10}\left[
 1+\frac{\lambda_1e^{-\lambda_2t}
 -\lambda_2e^{-\lambda_1t}}
 {\lambda_2-\lambda_1}\right]}.
\]
\excheck $N_3(0)=0$，而
$\lim_{t\to\infty}N_3(t)=N_{10}$，
符合所有母核最终变为稳定核的假定。

\radioactivityproblem{232}{258}
$\mathrm{^{210}Bi}\xrightarrow[\beta^-]{\lambda_1}
\mathrm{^{210}Po}\xrightarrow[\alpha]{\lambda_2}
\mathrm{^{206}Pb}$（稳定），
$\lambda_1=1.60\times10^{-6}\,\mathrm{s^{-1}}$、
$\lambda_2=5.80\times10^{-8}\,\mathrm{s^{-1}}$。
初始纯 $\mathrm{^{210}Bi}$ 的质量 $1.00\,\mathrm{mg}$。
求制成一个月后的 $\alpha$、$\beta$ 活度。
\exsolution
初始母核数
$N_0=(10^{-3}\,\mathrm g/210\,\mathrm{g/mol})N_A$。
取一月30天，$t=30(86400)\,\mathrm s$。
母核数 $N_{\rm Bi}=N_0e^{-\lambda_1t}$，
因此
\[
 \boxed{A_\beta=\lambda_1N_0e^{-\lambda_1t}
 \simeq7.25\times10^{10}\,\mathrm{Bq}}.
\]
子核
\[
 N_{\rm Po}=
 \frac{\lambda_1N_0}{\lambda_2-\lambda_1}
 (e^{-\lambda_1t}-e^{-\lambda_2t}),
\]
所以
\[
 \boxed{A_\alpha=\lambda_2 N_{\rm Po}
 \simeq1.46\times10^{11}\,\mathrm{Bq}}.
\]
\excheck 题目给出两种不同衰变的活度，
应分别使用瞬时母、子核数，
而不是将初始母核数直接乘两个常数。

\radioactivityproblem{233}{258}
(a) $\mathrm{^{226}Ra}$ 经五次 $\alpha$
和四次 $\beta^-$ 衰变生成何核？
(b) $\mathrm{^{238}U}$ 最终变为稳定
$\mathrm{^{206}Pb}$，经历多少次 $\alpha$
与 $\beta^-$ 衰变？
\exsolution
一次 $\alpha$ 衰变使
$(A,Z)\mapsto(A-4,Z-2)$；
一次 $\beta^-$ 使
$(A,Z)\mapsto(A,Z+1)$。
(a) Ra 初始 $(226,88)$，
五次 $\alpha$ 后 $(206,78)$，
四次 $\beta^-$ 后 $(206,82)$，
故 $\boxed{\mathrm{^{206}Pb}}$。
(b) 质量数从238到206，减少32，
需要 $\boxed{n_\alpha=32/4=8}$。
铀初始 $Z=92$，
八次 $\alpha$ 后 $Z=76$，
到铅$Z=82$还需
$\boxed{n_{\beta^-}=6}$。
\excheck $\beta^-$ 不改变质量数，
因此先由质量数差求 $\alpha$ 次数更直接。

\radioactivityproblem{234}{258}
假设初始静止的 $\mathrm{^{206}Pb}$ 核
发生 $\alpha$ 衰变，粒子动能
$T_\alpha=5.77\,\mathrm{MeV}$。
求剩余娘核的反冲速度，
以及反冲能占释放总能量的份额。
\exsolution
衰变前总动量为零，因此娘核动量与 $\alpha$
动量大小相等反向。
非相对论近似
$p_\alpha=\sqrt{2m_\alpha T_\alpha}$，
故娘核速率
\[
 \boxed{v_d=\frac{\sqrt{2m_\alpha T_\alpha}}{M_d}
 \simeq3.3\times10^5\,\mathrm{m/s}},
 \quad M_d\simeq202u.
\]
反冲能比
$T_d/T_\alpha=m_\alpha/M_d$，
故在总动能
$Q=T_\alpha+T_d$ 中，
\[
 \boxed{\frac{T_d}{Q}
 =\frac{m_\alpha}{M_d+m_\alpha}
 \simeq\frac4{206}=1.94\%\approx2.0\%}.
\]
\excheck 本题按原书是\emph{假想}的
$\mathrm{^{206}Pb}$ 衰变运动学模型，
并不意味着稳定 $\mathrm{^{206}Pb}$ 真实自发 $\alpha$ 衰变。

\radioactivityproblem{235}{258--259}
质量 $1.00\,\mathrm{mg}$ 的 $\mathrm{^{210}Po}$
制剂在其平均寿命时间内释放多少热？
每次 $\alpha$ 粒子动能为 $5.3\,\mathrm{MeV}$，
娘核几乎都形成基态。
\exsolution
平均寿命 $\tau=1/\lambda$ 时，
已衰变份额
$f=1-e^{-1}$。
核数
$N_0=(m/M)N_A$。
因娘核反冲，
总释放能量略大于 $\alpha$ 动能：
$Q=T_\alpha(1+m_\alpha/M_d)
\simeq T_\alpha(210/206)$。
于是
\[
 \boxed{W=N_0(1-e^{-1})Q
 \simeq1.57\times10^6\,\mathrm J
 \approx1.6\,\mathrm{MJ}}.
\]
\excheck 经过一个平均寿命仅约63.2\%的核衰变，
不能把所有核都当作已经衰变。

\radioactivityproblem{236}{259}
$\mathrm{^{210}Po}$ 的 $\alpha$ 衰变有两组
$\alpha$ 动能 $5.30$ 与 $4.50\,\mathrm{MeV}$；
前者使娘核至基态，后者至激发态。
求退激时 $\gamma$ 光子的能量。
\exsolution
母、娘核静质量差相同，
两衰变分支的总释放能之差
等于娘核激发能。
反冲动能比为$m_\alpha/M_d$，
所以
\[
 E^*=(T_{\alpha,1}-T_{\alpha,2})
 \left(1+\frac{m_\alpha}{M_d}\right).
\]
取$M_d\simeq206u$，
\[
 \boxed{E_\gamma\simeq E^*
 =(5.30-4.50)\frac{210}{206}
 \simeq0.816\,\mathrm{MeV}\approx0.82\,\mathrm{MeV}}.
\]
\excheck $\gamma$ 发射也有微小核反冲，
本题 MeV 精度下可以忽略。

\radioactivityproblem{237}{259}
常温常压空气中 $\alpha$ 粒子平均射程经验式
$R=0.98\times10^{-27}v_0^3\,\mathrm{cm}$，
其中 $v_0$ 的单位为 $\mathrm{cm/s}$。
初始动能 $T_\alpha=7.0\,\mathrm{MeV}$。
求 (a)射程；(b)全程及前半段射程内
平均产生的离子对数，每对形成能为 $34\,\mathrm{eV}$。
\exsolution
初速度
$v_0=\sqrt{2T_\alpha/m_\alpha}
\simeq1.84\times10^9\,\mathrm{cm/s}$，
代经验式
\[
 \boxed{R\simeq6.1\,\mathrm{cm}}.
\]
全部动能转为电离的近似下，
\[
 \boxed{N_{\rm total}=\frac{7.0\times10^6}{34}
 \simeq2.06\times10^5}.
\]
为计算前半段，用残余射程
$R_{\rm rem}\propto v^3\propto T^{3/2}$，
因此通过前半段距离后残余射程为$R/2$，
剩余动能
$T_{\rm mid}=T_0\,2^{-2/3}$。
前半段的能量损失为
$\Delta T=T_0(1-2^{-2/3})$，所以
\[
 \boxed{N_{\rm first\,half}=
 \frac{7.0\times10^6}{34}(1-2^{-2/3})
 \simeq7.62\times10^4}.
\]
\excheck 能量损失沿射程并不均匀，
故前半段产生的离子对不是总数的一半。
